%======================================================================
%  Supplementary Methods -- origin classification (paper section 3)
%
%  Source artifact: artifact_origin/page.html, section#methods,
%  snapshot 2026-08-28. Nine subsections, transcribed verbatim.
%  The artifact carries no display equations and no citation markers; every
%  named tool below has a "% TODO cite:" at its first mention in this file
%  or in 03_origin.tex.
%======================================================================

\subsection*{Reference origin curation}

\emph{oriT} references were compiled from three sources: the oriTDB2 \cite{liu2025oritdb} database, six consensus \emph{oriT} sequences from an \emph{oriT}-diversity
study,
% TODO cite: the oriT-diversity study. NOT resolved -- two candidates fit and
% attaching the wrong one would be a mis-citation, so this is left for the
% author. Candidate A: aresarroyo2023origins, whose Methods name six oriT
% clusters (F-like, R6K-like, R64-like, ColE1-like, RP4-like, R46-like) -- but
% that is the same paper as the transfer-network source below, and the text
% here presents them as two separate sources. Candidate B: zrimec2020multiple,
% which frames six MOB groups {F,P,Q,V,C,T} over 112 oriT regions. A third
% possibility is aresarroyo2024expanding. Pick the one you actually used.
and 87 experimentally validated \emph{oriT}s from an
\emph{oriT}-transfer-network study \cite{aresarroyo2023origins}.
% FLAG (2026-09-04): the COUNT does not match the cited paper. Ares-Arroyo
% 2023 states a final dataset of 91 origins of transfer, not 87, and no
% oriT-network study reporting ~87 could be found. Either correct 87 to 91, or
% name the source you took 87 from and this citation changes with it.
Experimental oriTDB2 entries were manually curated against literature
coordinates (for example oriT\_pIP421 extended from 7 to 194 bp spanning
IHF1--IR3, oriT\_CloDF13 from 13 to 661 bp), and two entries were
reclassified from bioinformatic to experimental. Bioinformatic-predicted
\emph{oriT}s were retained only if they were not exact sequence or identifier
duplicates of the experimental set and were $\geq$15 bp (2,860 sequences).
These 2,860 bioinformatic entries were not used as BLAST queries for
classifier labels. The final experimental set of 142 sequences was reduced to
\textbf{112} by MMseqs2 \cite{steinegger2017mmseqs2} clustering (minimum sequence identity 90\%, coverage 90\%, coverage-mode =
shorter). \emph{oriV} references (452 sequences from the OriV-Finder
database \cite{li2025orivfinder}) were deduplicated identically to \textbf{224} sequences.

\subsection*{Plasmid databases and origin detection}

Origins were located in two plasmid collections: PLSDB \cite{molano2025plsdb} (filtered to $\geq$ 2,000 bp) and IMG/PR \cite{camargo2024img}.
Representative selection used a 95\% ANI, 95\% unidirectional
alignment-fraction grouping (\texttt{ani95\_af95\_one\_direct}), which is
distinct from the 99\% ANI dereplication columns. Reference origins were
aligned to plasmid genomes with BLASTN \cite{camacho2009blast} (word size 7, dust filtering off),
pre-filtered at 50\% subject coverage and then retained at $\geq$ 80\%
identity and $\geq$ 80\% subject coverage; overlapping hits were merged and
the longest-span alignment kept per plasmid. After PLSDB selection, IMG/PR
candidates matching a selected PLSDB plasmid at ANI $\geq$95\% and
max(alignment fraction) $\geq$95\% (skani \cite{shaw2023fast}) were excluded.

\subsection*{Intergenic regions and plasmid selection}

% Pyrodigal-gv has no publication of its own; its README asks for these three.
Coding sequences were predicted with Pyrodigal-gv
\cite{larralde2022pyrodigal,hyatt2010prodigal,camargo2024genomad} in
metagenomic mode; intergenic regions were the complement of all CDS
intervals with a minimum length of 50 bp, and plasmids with fewer than three
qualifying intergenic regions were excluded. For each origin, overlap with
all intersecting intergenic regions was summed; the origin was retained if
that total was $\geq$ 50\% of the origin length, and every intersecting
region was then emitted. Only plasmids in which all detected origins met this
intergenic criterion were retained. Within each 95\% ANI grouping the
median-genome-length plasmid was taken as the representative, and for each
unique origin the top five plasmids by alignment quality (percent identity
$\times$ percent coverage) were kept; all origins of any selected plasmid
were retained to preserve multi-origin plasmids, and plasmids with more than
five origins were excluded as suspected assembly artefacts. For IMG/PR,
plasmids carrying frequently observed origins (occurrence above the mean + 1
s.d.) were required to be annotated as putatively complete, and for
\emph{oriT} to additionally carry mobilization genes; PLSDB plasmids were
assumed complete.

\subsection*{Homology-partitioned cross-validation splits}

To limit similarity leakage, parts were defined at the level of the origin
sequences, not the plasmids. All-against-all pairwise similarity was computed
with EMBOSS Water \cite{rice2000emboss}
(Smith--Waterman \cite{smith1981identification}, gap open 10, gap extend 0.5) on both the forward and
reverse-complement orientations, keeping the higher-scoring alignment;
distance was $1 - (\text{identity}/100)$, with pairs below 70\% identity or
50\% coverage set to the maximum distance. GraphPart \cite{teufel2023graphpart} partitioned the origins into five parts at a similarity threshold of 0.3 with
no post-hoc node moving during partitioning. For \emph{oriT}, GraphPart
retained 87 of 112 origins; the 25 removed origins were reinserted into the
part of their nearest assigned neighbour (ties to the smaller part), giving
final part sizes 26/23/23/21/19. After reinsertion, some cross-part pairs
fall inside the nominal 0.3 distance, so the \emph{oriT} split is
similarity-aware rather than strictly homology-separated. \emph{oriV}
retained all 224 origins with part sizes 46/46/44/44/44. Plasmids inherited
their part through their origins, and plasmids whose origins fell in
different parts were removed (retained selected plasmids: \emph{oriT} 381/399
PLSDB and 457/464 IMG/PR; \emph{oriV} 563/688 PLSDB and 582/686 IMG/PR). Each
cross-validation split trains on three parts, validates on one, and tests on
one (for split \emph{N}, part \emph{N}$-$1 tests and part \emph{N}
validates).

\subsection*{Windowing and class imbalance}

The window dataset combined PLSDB and IMG/PR at the intergenic-region,
origin-annotation and genome levels. Genomes were tiled into 512 bp windows
at a 256 bp step for training-set construction. Each window was assigned to
its maximum-overlap intergenic region; windows that overlapped only CDS were
discarded. A window was labelled positive when it overlapped an annotated
origin by $\geq$ 50\% of min(origin length, 512 bp) or by $\geq$ 256 bp.
Partial-origin windows below both thresholds were removed. Where a positive
intergenic region on a chromosome had no grid-positive window, a single
origin-centred window was inserted. Reverse-complement augmentation then
duplicated every retained window. The evaluated population is therefore this
filtered intergenic-window set, not every sliding window on the plasmid.
Region-level neg:pos was $\approx$ 50:1 (\emph{oriT}) and $\approx$ 20:1
(\emph{oriV}) before downsampling. Negatives were then selected
deterministically within each part: plasmids below the target ratio
contributed all negatives, and plasmids above it contributed nearest
negatives, ordered by distance to the closest positive window, in round-robin
until the part reached exactly 20:1. The reported training runs used a fixed
positive-class loss weight of 20$\times$ for both tasks.

\subsection*{Sequence embedders and probes}

Four families of frozen feature extractors were compared. Evo2 (base) at
three scales --- 1B, 7B and 40B parameters --- supplied embeddings from hook
\texttt{blocks.\{i\}.mlp.l3} over probed layer ranges of 10--24, 10--31 and
10--22 respectively; \model supplied embeddings from layers 36--72
(\texttt{blocks.\{i\}}); a non-pretrained one-hot nucleotide encoding served
as the from-scratch baseline; and PlasmidGPT (a BPE transformer pretrained on
engineered Addgene \cite{kamens2015addgene} plasmids) served as a negative control. PlasmidGPT token hidden states were
upsampled to nucleotide resolution by repeating each token vector across the
nucleotides spanned in the tokenizer offset mapping. Contextual embedders
received extra flanking bases in the forward pass (768 bp per side for
\model; 3,840 bp per side for Evo2 and PlasmidGPT, as set by the embedding
orchestration), discarded before saving, so each 512 bp window yielded
exactly 512 embedding vectors, with all layers of interest extracted in one
pass and cached. On the contextual embeddings a mean-pooled MLP probe emitted
one logit per window; on one-hot sequence a convolutional probe (CNN and
residual CNN; hidden 64, kernel 16, dropout 0.15) integrated context
spatially before emitting a logit.

\subsection*{Training and layer selection}

Probes were trained with five-fold cross-validation at the origin level
(GraphPart parts mapped to five plasmid parts; for split \emph{N}, part
\emph{N}$-$1 tests, part \emph{N} validates, the remaining three train), with
five stochastic runs per (embedder, layer, split) and no explicit Python,
NumPy or PyTorch seeds, giving 25 probes per (embedder, layer). The loss was
BCE-with-logits with the positive class up-weighted by a fixed factor of
20$\times$ for both tasks; a per-window loss-weight mechanism exists but was
disabled in the reported runs (all per-window weights were 1). Optimization
used AdamW \cite{loshchilov2019adamw} (learning rate $5\times10^{-4}$, weight decay 0.01) with cosine annealing
over 30 epochs and mixed precision, at batch size 256. Gradient accumulation
remained at its default of one and was not used as an active technique. The
best epoch per run was selected by validation AUPRC rather than validation
loss, and both the held-out test and validation logits at that epoch were
persisted so calibration could be fit after training. For each (embedder,
probe) the reported layer was the one maximizing validation AUPRC averaged
over both tasks, with each task's value itself the mean over the five splits.
Layer selection therefore uses only validation columns, but because parts
rotate, each part is a validation set in one split and a test set in another.
Reporting considered GPN layers 60--72 even though probes were trained on
36--72; the selected layers were \model \texttt{blocks.60}, Evo2-1B/7B
\texttt{blocks.15.mlp.l3}, Evo2-40B \texttt{blocks.16.mlp.l3}, and PlasmidGPT
\texttt{h.3}. All PlasmidGPT layers were swept and its best layer still
underperformed the from-scratch CNN.

\subsection*{Calibration and region aggregation}

Raw logits under class imbalance are poorly calibrated and not comparable
across runs, splits or embedders. Each probe's logits were mapped to
probabilities by isotonic regression \cite{zadrozny2002isotonic}
(\texttt{Isotonic\allowbreak Regression(\allowbreak out\_of\_\allowbreak bounds=\allowbreak"clip")})
fit on that probe's own
validation logits. For region-level inference, each strand's logits were
calibrated first and the corresponding forward and reverse-complement
probabilities were then averaged. The five per-run probability vectors per
(embedder, layer, split) were averaged by arithmetic-mean soft voting. For
region-level scoring, per-region window probabilities were aggregated with
\texttt{percentile-75} --- the 75th percentile when a region has more than
five windows, otherwise the maximum --- and a second isotonic calibrator was
fit per split on those aggregated validation scores. This aggregator is a
plain percentile and does not require adjacent windows or a contiguous span
of evidence. Extra inference span beyond the intergenic interval was
estimated as inference-region length minus intergenic length, not by direct
CDS intersection. After isotonic regression both the 75th and 95th
percentiles are well calibrated on the reliability diagrams; the 75th
percentile spreads probability mass more smoothly across the ranking range.

% NOTE (2026-09-04): the artifact titled this subsection simply "Inference",
% which collided with the pretraining Methods' own "Inference" heading. With
% both Methods blocks now in one main-article Methods section the collision was
% visible in the running text, so the two were disambiguated: the pretraining
% one is "Pretraining inference", this one is "Origin inference and deployment".
\subsection*{Origin inference and deployment}

At inference, intergenic regions were tiled into 512 bp windows at a 16 bp
stride, denser than the 256 bp training step used only for example
construction, and each region was embedded on both the forward and
reverse-complement strands. For the window-level benchmark, the forward and
reverse-complement of each window were treated as independent samples through
the probe and its per-run calibrator. For region-level discovery, the two
strands' calibrated probabilities were averaged so each genomic window
contributed one value to aggregation. Regions longer than the embedder's
maximum length were padded symmetrically to a multiple of that length, split
into non-overlapping chunks, embedded, concatenated, and trimmed back to the
original region; chunk order was reversed on the reverse-complement strand.
Context is added around the overall padded region rather than independently
around every chunk. Genome-wide application used standalone \model: five
calibrated runs within each of five splits (25 probes), per-split region
calibration, then an unweighted mean across splits. That mean is an ensemble
score, not a globally recalibrated probability. The Evo2-7B $+$ \model
ensemble was evaluated only as a benchmark combination.

Novelty protocol. Held-out positive windows were stratified by BLASTN
similarity to that split's training positives. Similarity is percent identity
$\times$ query coverage / 100; windows with score $\leq$20 are novel and
those with score $>$20 are training-similar. Forward and reverse-complement
probabilities were averaged before scoring. All positives were retained;
negatives were sampled without replacement at 20:1 over 200 deterministic
draws; average precision was computed per split and macro-averaged. The BLAST
tables used for this stratification are not regenerated in the current
repository, so the claim is limited to the operational score threshold rather
than to the absence of any homolog.

% Added 2026-09-04. Both paragraphs carry facts that previously appeared ONLY
% in the Results prose, which was trimmed to the Nature Biotechnology main-text
% limit. Nothing here is new: the timing harness and the reference-tool
% database ablation are transcribed from 03_origin.tex as they stood.
\subsection*{Reference-tool baselines}

The reference tools were run as region rankers over the same intergenic
regions supplied to the probes, with oriTfinder2 serving the \emph{oriT} task
and OriV-Finder the \emph{oriV} task. Because these tools match a query
against a curated database rather than learning from a training fold, their
exposure to the benchmark is controlled by which origins that database is
allowed to contain. Restricting each tool's database to the training fold
alone, rather than to the training and validation folds together, changed its
mean score by at most 0.007 on either task, so the reported comparison is not
sensitive to that choice. Individual splits move further, by up to 0.024 for
\emph{oriT}.
% Corrected 2026-09-04: the bound was stated without saying it applies to the
% mean. Per-split differences reach 0.024 (oriT split 0) and 0.008 (oriV split
% 4), both above the stated bound, and the supplementary figure this cites is a
% per-split heatmap, so a reader checks it per split.

\subsection*{Inference-cost measurement}

Per-region inference cost was measured with a single-split, five-run window
predictor over 794 intergenic regions drawn from 10 plasmids, timed on one
CUDA device. The comparison holds the region set, the window size, the stride
and the number of stochastic runs fixed across embedders, so the reported
difference reflects the embedder alone and not the surrounding pipeline.
